\def\writeANS{}
\begin{loigiaibth}{7}
  \begin {eqnarray} \Delta _r H_{298}^\circ &=& 4.\Delta _f H_{298}^\circ (CO_2) + 2.\Delta _f H_{298}^\circ (H_2O) - 2.\Delta _f H_{298}^\circ (C_2H_2) - 5.\Delta _f H_{298}^\circ (O_2) \nonumber \\ -2243{,}6 &=& 4.(-393{,}5) + 2.(-285{,}8) - 2.\Delta _f H_{298}^\circ (C_2H_2) - 5.0 \nonumber \\ &\Rightarrow & \Delta _f H_{298}^\circ (C_2H_2) = +49 \text { kJ/mol} \nonumber \end {eqnarray}
\end{loigiaibth}
\def\writeANS{}
\begin{loigiaibth}{8}
  \begin {enumerate}[(1)] \item $\Delta _r H_{298}^\circ (1) = -1206{,}9 - (-635{,}1) - (-393{,}5) = -178{,}3 kJ \approx -178 \ kJ$. \item $\Delta _r H_{298}^\circ (2) = -393{,}5 - 0 - 0 = -393{,}5 kJ \approx -394 \ kJ$. \item $\Delta _r H_{298}^\circ (3) = [\frac {1}{2}.(-393{,}5) + (-269{,}8)] - [\frac {1}{2}.87{,}9 - \frac {3}{2}.0] = -510{,}5 \ kJ \approx -511 \ kJ$ \item $\Delta _r H_{298}^\circ (4) = [1.0 + 3.(-241{,}82)] - [2.(-45{,}9) + \frac {3}{2}.0] = -633{,}66 \ kJ \approx -634 \ kJ$ \end {enumerate}
\end{loigiaibth}
\def\writeANS{}
\begin{loigiaibth}{9}
 Enthalpy của quá trình: $\Delta _r H_{298}^\circ = -542{,}83 + 2.(-167{,}16) - (-795{,}0) = -82{,}15 \ kJ \approx -82 \ kJ$.
\end{loigiaibth}
\def\writeANS{}
\begin{loigiaibth}{10}
 $\Delta _r H_{298}^\circ = 6 \times (-393{,}50) + 5 \times (-241{,}82) - 2 \times (-370{,}15) = -2829{,}8 \ kJ$ \par Nhiệt tỏa ra khi tạo thành 1 mol $CO_2$ là $\frac {2829{,}8}{6} = 471{,}63 \approx 472 \ kJ$.
\end{loigiaibth}
\def\writeANS{}
\begin{loigiaibth}{11}
 $Fe_2O_3(s) + 3CO(g) \longrightarrow 2Fe(s) + 3CO_2(g)$ \par $\Delta _r H_{298}^\circ = 3 \times (-393{,}5) - 2.0 - 3.(-110{,}5) - (-824{,}4) = -24{,}8 \ kJ$ \par Khi cho 1 mol $Fe_2O_3$ pư với 1 mol CO thì CO hết $\Rightarrow $ nhiệt lượng giải phóng = $24{,}8/3 = 8{,}27 \ kJ$.
\end{loigiaibth}
\def\writeANS{}
\begin{loigiaibth}{12}
 \begin {eqnarray} \Delta _r H_{298}^\circ (1) &=& -162{,}8 \text { kJ} \nonumber \\ \Delta _r H_{298}^\circ (2) &=& -50{,}1 \text { kJ} \nonumber \\ \Delta _r H_{298}^\circ (3) &=& -73{,}3 \text { kJ} \nonumber \\ \Delta _r H_{298}^\circ (4) &=& -45{,}6 \text { kJ} \nonumber \\ \Delta _r H_{298}^\circ (5) &=& -66{,}4 \text { kJ} \nonumber \end {eqnarray}
\end{loigiaibth}
\def\writeANS{}
\begin{loigiaibth}{13}
 Lượng nhiệt cần để thu được $0{,}1$ mol CaO là $0{,}1 \times 178{,}49 = +17{,}849 \ kJ$. \par (a) Lượng $C_2H_5OH(l)$ cần dùng: $\frac {17{,}849}{1370{,}7} = 0{,}013 \ mol \Rightarrow 0{,}6 \ g$. \par (b) Lượng $C(graphite,s)$ cần dùng: $\frac {17{,}849}{393{,}509} = 0{,}045 \ mol \Rightarrow 0{,}54 \ g$.
\end{loigiaibth}
